\section{이론적 분석}
\label{sec:이론적분석}

본 장에서는 제안하는 글로벌-로컬 정책 증류 기법의 이론적 성능 한계를 도출하고, 예측적 위험 스위칭이 부분 관측성 하에서 라우팅 오류를 어떻게 제어하는지 분석한다.

\subsection{로컬 배포 가능성과 계산 복잡도}
분산형 UAV 라우팅 시스템이 성립하기 위해서는 최적의 포워딩 결정 $a^\star$이 전역 토폴로지 정보 $\mathcal{G}_t$나 타 무인기의 글로벌 탐색 상태 없이 오직 자신의 1-hop 정보 $o_u(t)$만을 기반으로 도출되어야 한다. \method가 제안하는 모든 특징량($m_i, \ell_i, q_i, o_i, \bar{o}_{i,\mathrm{top}k}, \rho_i, p_{i,\mathrm{keep}}, \eta_i$)은 1-hop 로컬 주기적 비콘 메시지 교환 및 수신 전력 판정만을 사용하므로 온보드 분산 실행이 가능하다. 최종 결정을 위한 시간 복잡도는 이웃 노드 수에 선형적인 $\mathcal{O}(|\mathcal{N}_u(t)|)$를 유지하며, 이는 자원이 한정된 UAV 온보드 비행 제어 컴퓨터에서 실시간으로 완수될 수 있는 수준이다.

\subsection{라우팅 오차 분해 이론}
전역 그래프 정보를 관측하는 오프라인 교사 정책 $\pi_T(a \mid s)$와 로컬 정보만을 관측하는 학생 정책 $\pi_S(a \mid o)$ 간의 총 라우팅 성능 차이 $\mathcal{E}_{\mathrm{total}}$은 정보 불일치에 의한 오차와 근사 오차의 합으로 공식 분해할 수 있다.
\begin{equation}
    \mathcal{E}_{\mathrm{total}} = \mathcal{E}_{\mathrm{info}} + \mathcal{E}_{\mathrm{approx}}
\end{equation}
\begin{enumerate}
    \item \textbf{정보 격차 오차(Information-Gap Error)} $\mathcal{E}_{\mathrm{info}}$: 로컬 매핑 함수 $o = \mathcal{O}(s, u)$의 비단사(Non-injective) 특성 때문에 발생한다. 즉, 여러 다른 전역 토폴로지가 동일한 로컬 관측 상태로 매핑되는 관측 칭의성(Observation Aliasing)으로 인해 극복 불가능한 정보 소실이 발생한다.
    \item \textbf{근사 및 최적화 오차(Approximation/Optimization Error)} $\mathcal{E}_{\mathrm{approx}}$: 제한된 노드 파라미터 용량 및 오프라인 경사하강법 최적화의 한계로 인해 발생하는 오차이다.
\end{enumerate}

\subsection{성능 바운드 및 위험 스위칭의 오차 제어 메커니즘}
Performance Difference Lemma에 의하면, 교사 모델 대비 학생 모델의 리턴 성능 격차 $\Delta J = J(\pi_T) - J(\pi_S)$는 학생 정책의 상태 점유 분포(occupancy distribution) $d^{\pi_S}$ 상에서의 평균 KL divergence로 상한선(Bound)이 정의된다.
\begin{equation}
    \Delta J \le \frac{1}{1-\gamma} \mathbb{E}_{s \sim d^{\pi_S}} \left[ \sqrt{2 D_{\mathrm{KL}}(\pi_T(\cdot \mid s) \Vert \pi_S(\cdot \mid \mathcal{O}(s)))} \right]
\end{equation}
일반적인 지식 증류 환경에서는 평균적인 상태 분포 상의 KL divergence만을 최소화하므로, 라우팅 홀이나 단절 직전의 링크 같은 드물게 마주하는 예외 상태(Out-of-Distribution, OOD)에서의 대처 능력이 상실된다. 학생 정책이 이러한 임계 상태에서 잘못된 결정을 내리면 네트워크 전달 실패 및 루핑 등 치명적인 성능 저하(Return Gap의 급격한 증가)가 초래된다.

제안하는 예측적 위험 스위칭(PRS) 모듈은 이러한 리턴 갭 상한선 증가를 동적으로 제어한다. 평상시의 상태에서는 nominal branch가 수행되어 스위치 활성 확률이 영에 수렴한다.
\begin{equation}
    \Pr(S=1 \mid D_{a_N} \le \tau) \approx 0
\end{equation}
그러나 위험 요인이 감지되어 nominal 경로의 danger score가 임계치 $\tau$를 넘는 순간, 스위치가 활성화되어 강인한 예측 노드로 차홉을 즉시 편향시킨다. 이를 통해 OOD 위험 상태가 점유 분포 $d^{\pi_S}$ 내에서 장기적으로 유지되거나 누적되는 것을 방지함으로써, 종단간 성능 격차 $\Delta J$를 안정적으로 제어할 수 있게 된다.
